\section{Урок 7: интерактивное управление и запуск Lua-сценариев}

\subsection{Цели урока}
\begin{itemize}
\item Понять, как устроен ручной канал управления через клавиатуру.
\item Освоить отправку RC-каналов и логику привязки клавиш.
\item Научиться запускать и останавливать Lua-сценарий с Python.
\end{itemize}

\subsection{Опорные листинги}
\lstinputlisting[style=GeoskanStyle,language=pythoncustom,caption={Ручное управление через WASD и RC-каналы}]{code/python/examples/WASD_flight.py}
\lstinputlisting[style=GeoskanStyle,language=pythoncustom,caption={Управление Lua-скриптом на борту}]{code/python/examples/LUA_script_start.py}

\subsection{Ключевые команды}
\begin{itemize}
\item \texttt{send\_rc\_channels(channel\_1,\dots,channel\_5)}
\item \texttt{lua\_script\_control("Start")} и \texttt{lua\_script\_control("Stop")}
\item \texttt{cv2.waitKey(1)}
\end{itemize}

\begin{warnbox}
Интерактивные скрипты подходят для тренировки пилотирования и диагностики, но не заменяют автономные безопасные миссии.
\end{warnbox}

\begin{taskbox}[Минимальный набор упражнений]
\begin{enumerate}
\item Измените шаг каналов \texttt{min\_v/max\_v} и проверьте чувствительность.
\item Добавьте клавишу экстренной остановки с \texttt{land()}.
\item Реализуйте текстовую подсказку активной команды в консоли.
\end{enumerate}
\end{taskbox}

\subsection{Критерии успешности}
\begin{enumerate}
\item Все горячие клавиши вызывают ожидаемую реакцию.
\item При выходе через \texttt{ESC} соединение закрывается корректно.
\item Lua-скрипт запускается и останавливается по командам без перезапуска Python.
\end{enumerate}
